\subsection{交换环上的形式幂级数}

对每个集合 I，加法幺半群 \(\mathbf{N}^{(I)}\) 满足第 10 号中的条件 (D)；因为，若 \(s = (n_i)_{i \in I}\) ，其中 \(n_i = 0\) ，除非指标 \(i\) 属于 I 的某个有限子集 \(\mathbf{H}\) ，则关系 \(s = t + u\) （其中 \(t = (p_i)_{i \in I}\) ， \(u = (q_i)_{i \in I}\) ）等价于 \(p_i + q_i = n_i\) 对一切 \(i\) 成立；而这蕴含 \(p_i = q_i = 0\) （对 \(i \notin \mathbf{H}\) ）以及 \(p_i \leqslant n_i\) , \(q_i \leqslant n_i\) （对 \(i \in \mathbf{H}\) ）；因此有 \(\prod_{i \in \mathbf{H}} (n_i + 1)\) 个有序对 \((t, u)\) 在 \(\mathbf{N}^{(I)}\) 中满足 \(t + u = s\) .

因此，我们可以考虑幺半群 \(\mathbf{N}^{(I)}\) 在 A 上的全代数，它包含该幺半群的（限制）代数 \(\mathrm{A}[\mathbf{X}_i]_{i\in I}\) 。这是一个含单位元的交换结合代数，称为未定元 \(\mathbf{X}_i\) ( \(i\in I\) ） 、系数在 A 中的形式幂级数代数，记作 \(\mathrm{A}[[\mathbf{X}_i]]_{i\in I}\) ；其元素称为未定元 \(\mathbf{X}_i\) ( \(i\in I\) ） 、系数在 A 中的形式幂级数。这样的元素 \((\alpha_{\nu})_{\nu \in \mathbf{N}^{(I)}}\) 依照第 10 号中的约定也记作 \(\sum_{\nu \in \mathbf{N}^{(I)}}\alpha_{\nu}\mathbf{X}^{\nu}\) ； \(\alpha_{\nu}\) 称为该形式幂级数的系数， \(\alpha_{\nu}\mathbf{X}^{\nu}\) 称为它的项；因此， \(\mathbf{X}_i\) 的多项式就是只有有限多个非零（ \(\neq 0\) ）项的形式幂级数 .

显然，一个代数同构 \(\mathbf{A}[[\mathbf{X}_i]]_{i\in I_1}\to \mathbf{A}[[\mathbf{X}_i]]_{i\in I_2}\) 由每个双射 \(\sigma \colon I_1\to I_2\) 典范地导出，其办法是把形式幂级数 \(\sum_{(n_i)}\alpha_{(n_i)}\cdot \prod_{i\in I_1}\mathbf{X}_i^{n_i}\) 映为形式幂级数 \(\sum_{(n_i)}\alpha_{(n_i)}\cdot \prod_{i\in I_1}\mathbf{X}_{\sigma (i)}^{n_i}\) .

设 \(J\) 是 \(I\) 的子集；代数 \(A[[X_i]]_{i \in J}\) 可等同于 \(A[[X_i]]_{i \in I}\) 的一个子代数，它由满足下列条件的形式幂级数 \(\sum_{(n_i)} \alpha_{(n_i)} \cdot \prod_{i \in I} X_i^{n_i}\) 组成： \(\alpha_{(n_i)} = 0\) 对每个元素 \((n_i) \in \mathbf{N}^{(I)}\) （对该元素， \(n_i \neq 0\) 至少对一个指标 \(i \in I - J\) 成立）均成立。进一步，若 \(K = I - J\) ，则 \(A[[X_i]]_{i \in I}\) 典范地等同于 \((A[[X_j]]_{j \in J})[[X_k]]_{k \in K}\) ，其办法是把形式幂级数 \(\sum_{(n_i)} \alpha_{(n_i)} \cdot \prod_{i \in I} X_i^{n_i}\) 等同于形式幂级数 \(\sum_{(m_k)} \beta_{(m_k)} \cdot \prod_{k \in K} X_k^{m_k}\) ，其中

\[
\beta_ {(m _ {k})} = \sum_ {(p _ {j})} \gamma_ {(p _ {j})} \cdot \prod_ {j \in J} X _ {j} ^ {p _ {j}}
\]

其中 \(\gamma_{(p_i)} = \alpha_{(n_i)}\) 对这样的序列 \((n_i)\) 成立： \(n_i = p_i\) （对 \(i \in \mathbf{J}\) ）而 \(n_i = m_i\) （对 \(i \in \mathbf{K}\) ）.

给定形式幂级数 \(u = \sum_{\nu} \alpha_{\nu} X^{\nu}\) ，项 \(\alpha_{\nu} X^{\nu}\) （其中 \(|\nu| = p\) ）称为 \(u\) 中总次数 \(p\) 的项。形式幂级数 \(u_p\) （其总次数 \(p\) 的项即 \(u\) 的相应项，其余项为零）称为 \(u\) 的 \(p\) 次齐次部分；当 I 为有限集时， \(u_p\) 对一切 \(p\) 都是多项式； \(u_0\) 等同于 A 中的一个元素（也称为 \(u\) 的常数项）。若 \(u\) 与 \(v\) 是两个形式幂级数且 \(w = uv\) ，则

\[
w _ {p} = \sum_ {r = 0} ^ {p} u _ {r} v _ {p - r}\tag{39}
\]

对于每个整数 \(p \geqslant 0\) .

对每个形式幂级数 \(u \neq 0\) ，最小的整数 \(p \geqslant 0\) ，即使 \(u_{p} \neq 0\) 的最小者，称为 u 的全阶（或简称 u 的阶）。若这个阶记作 \(\omega(u)\) ，且 u 与 v 是两个 \(\neq 0\) 的形式幂级数，则

\[
\omega (u + v) \geqslant \inf (\omega (u), \omega (v)) \quad \text { if } u + v \neq 0\tag{40}
\]

\[
\omega (u v) \geqslant \omega (u) + \omega (v) \quad \text { if } u v \neq 0.\tag{41}
\]

进一步，如果 \(\omega(u) \neq \omega(v)\) ，那么必然 \(u + v \neq 0\) 且 (40) 的两边相等。

注意，0 的阶没有定义。按一种语言上的滥用，约定称“f 是一个阶 \(\geqslant p\) （相应地 \textgreater p ）的形式幂级数”，如果 f 的 n 次齐次部分对所有满足 n \textless{} p 的 n （相应地， \(n \leqslant p\) ）均为零；因此 0 是“阶 \textgreater p 的形式幂级数”，对每个整数 \(p \geqslant 0\) .

设 \(J\) 是 \(I\) 的子集，并设 \(A[[X_i]]_{i \in I}\) 如上所述地等同于 \(B[[X_k]]_{k \in K}\) ，其中 \(K = I - J\) ， \(B = A[[X_j]]_{j \in J}\) ；把上述定义应用于 \(B[[X_k]]_{k \in K}\) ，便对形式幂级数 \(u \in A[[X_i]]_{i \in I}\) 得到一些新的定义：项 \(\alpha_{(n_i)} \cdot \prod_{i \in I} X_i^{n_i}\) 称为次数为 \(p\) 的项（次数相对于 \(X_i\) 、指标 \(i \in K\) 而言），若 \(\sum_{i \in K} n_i = p\) ；而 \(B[[X_k]]_{k \in K}\) 中，其 \(p\) 次项与 \(u\) 的相同、其余项为零的形式幂级数，称为 \(p\) 次齐次部分（相对于 \(X_i\) 、指标 \(i \in K\) ）。若 \(u \neq 0\) ，阶 \(\omega_K(u)\) （相对于 \(X_i\) 、指标 \(i \in K\) ）是最小的整数 \(p \geqslant 0\) ，使得 \(u\) 的 \(p\) 次齐次部分（相对于 \(X_i\) 、指标 \(i \in K\) ）为 \(\neq 0\) 。当 \(\omega\) 换成 \(\omega_K\) 时，不等式 (40) 与 (41) 仍然成立 .

